% META: {"id": "TCOMPL-ME-CO-001", "chapitre": "TCOMPL-MODELES-EVOLUTION", "type_objet": "corrige", "exercice_ref": "TCOMPL-ME-EX-001", "capacites_codes": ["C2"], "sources_inspiration": [], "status": "generated"}
% BEGIN-VERIFY
% from sympy import Rational
% u0, q = 200, Rational(1, 4)
% limite = u0 / (1 - q)
% assert limite == Rational(800, 3)
% # La nouvelle question 1 demande la STRUCTURE : la somme partielle
% # doit valoir la formule de la somme geometrique, pas seulement sa limite.
% from sympy import symbols as _sym, summation as _sum, simplify as _simp
% _k, _n = _sym('k n', integer=True, nonnegative=True)
% _S = _sum(200 * Rational(1, 4) ** _k, (_k, 0, _n))
% assert _simp(_S - 200 * (1 - Rational(1, 4) ** (_n + 1)) / (1 - Rational(1, 4))) == 0
% END-VERIFY
\begin{corrige}{TCOMPL-ME-EX-001}
\textbf{1.} L'apport du jour $0$ vaut $200$ kg ; au bout de $k$ jours il n'en
subsiste que $200\times0{,}25^{\,k}$. En sommant les apports successifs, la
quantité présente au bout de $n$ jours vaut
\[
  q_n = \sum_{k=0}^{n} 200\times0{,}25^{\,k},
\]
somme des termes d'une suite géométrique de premier terme $200$ et de raison
$0{,}25$.

\textbf{2.} La raison $0{,}25$ appartient à $\left]0\,;\,1\right[$, donc
$0{,}25^{\,n+1}\to0$ et la somme converge :
\[
  q_n = 200\times\frac{1-0{,}25^{\,n+1}}{1-0{,}25}
  \xrightarrow[n\to+\infty]{} \frac{200}{0{,}75}
  = \frac{800}{3} \approx 266{,}7 \text{ kg}.
\]
\end{corrige}
